\documentclass[12pt,a4paper]{article}
\usepackage[utf8]{inputenc}
\usepackage[brazil]{babel}
\usepackage[T1]{fontenc}
\usepackage{amsmath}
\usepackage{amsfonts}
\usepackage{amssymb}
\usepackage{graphicx}
\usepackage{booktabs}
\usepackage{array}
\usepackage{geometry}
\usepackage{hyperref}
\usepackage{float}
\usepackage{xcolor}
\usepackage{multirow}
\usepackage{subcaption}
\usepackage{listings}
\usepackage{helvet}
\renewcommand{\familydefault}{\sfdefault}

\geometry{margin=2.5cm}
\graphicspath{{./imagens/}}

% Remover numeração das seções
\setcounter{secnumdepth}{0}

% Configuração para código Python
\lstset{
    language=Python,
    basicstyle=\ttfamily\small,
    keywordstyle=\color{blue},
    stringstyle=\color{orange},
    commentstyle=\color{gray},
    numbers=none,
    frame=single,
    breaklines=true,
    backgroundcolor=\color{gray!10},
    showstringspaces=false
}

\hypersetup{
    colorlinks=true,
    linkcolor=blue,
    filecolor=magenta,      
    urlcolor=cyan,
}

\begin{document}

% ============================================================================
% INTRODUÇÃO
% ============================================================================
\section{}

\textbf{Aluno:} Gabriel Francelino Voidaleski

\textbf{Curso:} Tecnologia em Análise e Desenvolvimento de Sistemas (TADS)

\textbf{Disciplina:} Tópicos Especiais em Banco de Dados 1

\vspace{0.5cm}

Este trabalho apresenta a aplicação de técnicas de aprendizado de máquina supervisionado para classificação de vinhos utilizando o dataset Wine da UCI. Foram implementados três algoritmos de classificação: Árvore de Decisão, Random Forest e SVM (Support Vector Machine). O objetivo é treinar os modelos, avaliar seu desempenho através de matrizes de confusão e métricas (acurácia, precisão, sensibilidade e F1-score), e identificar técnicas para melhorar os resultados.

% ============================================================================
% 2. BASE DE DADOS
% ============================================================================
\section{Base de Dados Escolhida}

\textbf{Dataset:} Wine Recognition Dataset

\textbf{Fonte:} UCI Machine Learning Repository

\textbf{URL:} \url{https://archive.ics.uci.edu/ml/datasets/wine}

\textbf{Descrição:} Dataset com análises químicas de vinhos de 3 cultivares diferentes da Itália. Contém 178 amostras, 13 atributos numéricos e 3 classes.

\begin{lstlisting}
# Carregar dados
from sklearn.datasets import load_wine
wine = load_wine()
X, y = wine.data, wine.target
\end{lstlisting}

% ============================================================================
% 2. TÉCNICAS DE APRENDIZADO
% ============================================================================
\section{Técnicas de Aprendizado Escolhidas}

\begin{itemize}
    \item \textbf{Árvore de Decisão} - Algoritmo simples e interpretável
    \item \textbf{Random Forest} - Ensemble de múltiplas árvores
    \item \textbf{SVM (Support Vector Machine)} - Classificador baseado em hiperplanos
\end{itemize}

\begin{lstlisting}
# Dividir dados (70% treino, 30% teste)
X_train, X_test, y_train, y_test = train_test_split(
    X, y, test_size=0.3, random_state=42)

# Treinar modelo inicial
modelo_inicial = DecisionTreeClassifier(random_state=42)
modelo_inicial.fit(X_train, y_train)
y_pred_inicial = modelo_inicial.predict(X_test)
\end{lstlisting}

% ============================================================================
% 3. MATRIZ DE CONFUSÃO - MODELO INICIAL
% ============================================================================
\section{Matriz de Confusão - Modelo Inicial}

O primeiro modelo treinado foi a Árvore de Decisão, utilizando os dados sem normalização. Este serve como baseline para comparação com os modelos otimizados.

\begin{figure}[H]
\centering
\includegraphics[width=0.6\textwidth]{matriz_inicial.png}
\caption{Matriz de Confusão - Árvore de Decisão (Modelo Inicial)}
\end{figure}

\begin{table}[H]
\centering
\caption{Valores da Matriz de Confusão - Árvore de Decisão}
\begin{tabular}{cc|ccc}
\toprule
& & \multicolumn{3}{c}{\textbf{Predito}} \\
& & Classe 0 & Classe 1 & Classe 2 \\
\midrule
\multirow{3}{*}{\textbf{Real}} 
& Classe 0 & 18 & 1 & 0 \\
& Classe 1 & 0 & 21 & 0 \\
& Classe 2 & 1 & 0 & 13 \\
\bottomrule
\end{tabular}
\end{table}

\textbf{Métricas do Modelo Inicial:}
\begin{itemize}
    \item Taxa de Acertos (Acurácia): 96,30\%
    \item Precisão: 96,38\%
    \item Sensibilidade (Recall): 96,30\%
    \item F1-Score: 96,28\%
\end{itemize}

% ============================================================================
% 4. MÉTODOS PARA MELHORAR
% ============================================================================
\section{Métodos Aplicados para Melhorar os Resultados}

Embora a Árvore de Decisão tenha alcançado 96,30\% de acurácia, buscamos melhorar os resultados através das seguintes técnicas:

\begin{itemize}
    \item \textbf{Normalização dos dados} (StandardScaler) - Padroniza os atributos para média 0 e desvio padrão 1
    \item \textbf{Troca de algoritmo} para Random Forest e SVM
\end{itemize}

\begin{lstlisting}
# Normalizar dados
scaler = StandardScaler()
X_train_scaled = scaler.fit_transform(X_train)
X_test_scaled = scaler.transform(X_test)

# Treinar Random Forest
modelo_rf = RandomForestClassifier(n_estimators=100, random_state=42)
modelo_rf.fit(X_train_scaled, y_train)

# Treinar SVM
modelo_svm = SVC(kernel='rbf', random_state=42)
modelo_svm.fit(X_train_scaled, y_train)
\end{lstlisting}

% ============================================================================
% 5. MATRIZ DE CONFUSÃO - RESULTADO OBTIDO
% ============================================================================
\section{Matriz de Confusão - Resultado Obtido}

\begin{figure}[H]
\centering
\includegraphics[width=0.95\textwidth]{matriz_rf_svm.png}
\caption{Matrizes de Confusão - Random Forest e SVM (Modelos Otimizados)}
\end{figure}

\subsection{Random Forest}

\begin{table}[H]
\centering
\caption{Valores da Matriz de Confusão - Random Forest}
\begin{tabular}{cc|ccc}
\toprule
& & \multicolumn{3}{c}{\textbf{Predito}} \\
& & Classe 0 & Classe 1 & Classe 2 \\
\midrule
\multirow{3}{*}{\textbf{Real}} 
& Classe 0 & 19 & 0 & 0 \\
& Classe 1 & 0 & 21 & 0 \\
& Classe 2 & 0 & 0 & 14 \\
\bottomrule
\end{tabular}
\end{table}

\textbf{Métricas Random Forest:}
\begin{itemize}
    \item Taxa de Acertos: 100,00\%
    \item Precisão: 100,00\%
    \item Sensibilidade: 100,00\%
    \item F1-Score: 100,00\%
\end{itemize}

\subsection{SVM}

\begin{table}[H]
\centering
\caption{Valores da Matriz de Confusão - SVM}
\begin{tabular}{cc|ccc}
\toprule
& & \multicolumn{3}{c}{\textbf{Predito}} \\
& & Classe 0 & Classe 1 & Classe 2 \\
\midrule
\multirow{3}{*}{\textbf{Real}} 
& Classe 0 & 19 & 0 & 0 \\
& Classe 1 & 0 & 21 & 0 \\
& Classe 2 & 0 & 1 & 13 \\
\bottomrule
\end{tabular}
\end{table}

\textbf{Métricas SVM:}
\begin{itemize}
    \item Taxa de Acertos: 98,15\%
    \item Precisão: 98,23\%
    \item Sensibilidade: 98,15\%
    \item F1-Score: 98,14\%
\end{itemize}

% ============================================================================
% 6. COMPARAÇÃO
% ============================================================================
\section{Comparação dos Resultados}

\begin{table}[H]
\centering
\caption{Resumo Comparativo}
\begin{tabular}{lccc}
\toprule
\textbf{Métrica} & \textbf{Decision Tree} & \textbf{Random Forest} & \textbf{SVM} \\
\midrule
Taxa de Acertos & 96,30\% & 100,00\% & 98,15\% \\
Precisão & 96,38\% & 100,00\% & 98,23\% \\
Sensibilidade & 96,30\% & 100,00\% & 98,15\% \\
F1-Score & 96,28\% & 100,00\% & 98,14\% \\
\bottomrule
\end{tabular}
\end{table}

\textbf{Melhor modelo:} Random Forest (100,00\%)

\textbf{Melhoria sobre inicial:} +3,70 pontos percentuais

% ============================================================================
% 7. CONCLUSÃO
% ============================================================================
\section{Conclusão}

Neste trabalho foram implementadas três técnicas de aprendizado de máquina supervisionado para classificação de vinhos: Árvore de Decisão, Random Forest e SVM. 

O modelo inicial (Árvore de Decisão sem normalização) obteve 96,30\% de acurácia. Após aplicar normalização dos dados com StandardScaler e testar algoritmos mais robustos, obtivemos melhorias significativas.

O Random Forest atingiu 100\% de acurácia no conjunto de teste, seguido pelo SVM com 98,15\%. Isso demonstra que a combinação de pré-processamento adequado (normalização) com algoritmos de ensemble (Random Forest) pode melhorar significativamente o desempenho em problemas de classificação.

\end{document}
